\hypertarget{appendix-0b3be18411f9}{}
\section*{Eigenvalues and eigenspaces}
\section{Eigenvalues and eigenspaces}\label{eigenvalues-and-eigenspaces}

An eigenvector is a nonzero direction that a square linear map scales by its eigenvalue. The eigenspace is the null space of \(A-\lambda I\). Real symmetric matrices have an orthonormal eigenbasis; general matrices may be defective or have sensitive eigenvectors.

\textbf{Use it when:} studying repeated powers, invariant directions, covariance, stability, or matrix functions.

\textbf{Developed in Phase 8:} Eigenanalysis begins with invariant directions, then distinguishes algebraic multiplicity from eigenspace dimension and shows why a defective matrix may not have a full eigenbasis. Power and inverse iteration target selected eigenpairs; QR iteration approaches a Schur form, while residual size and eigenvector sensitivity answer different accuracy questions. \href{Continuous\%20Notes/Computing\%20Foundations\%20-\%20From\%20Bits\%20to\%20C++\%20and\%20Systems\%20APIs\#Eigenvalues,\%20eigenspaces,\%20and\%20the\%20two\%20kinds\%20of\%20multiplicity}{eigenspaces and multiplicity} · \href{Continuous\%20Notes/Computing\%20Foundations\%20-\%20From\%20Bits\%20to\%20C++\%20and\%20Systems\%20APIs\#QR\%20iteration\%20and\%20Jacobi\%20rotations\%20preserve\%20the\%20spectrum\%20while\%20simplifying\%20shape}{eigenvalue iteration and Schur form}

\subsection{Prerequisites}\label{prerequisites}

\href{}{Matrices as linear maps and basis images (01-Foundations-I, appendix-bd39c3bf5d64)}, \hyperref[appendix-8572a4a060c1]{Linear systems and fundamental subspaces}, \hyperref[appendix-175ca5ff874d]{Matrix transpose symmetry and trace}. - \hyperref[appendix-ab31ef447dbf]{Matrix multiplication and composition order}

\subsection{Used by}\label{used-by}

\hyperref[appendix-1fad5e7f36fa]{Singular value decomposition}, \hyperref[appendix-ad175ef3c5bf]{PCA and covariance eigensystems}, \hyperref[appendix-73f5bcf7a462]{Polar decomposition}, \hyperref[appendix-0f4bc2bf2b1f]{Matrix functions}.

\textbf{Teaching path:} \href{Continuous\%20Notes/Computing\%20Foundations\%20-\%20From\%20Bits\%20to\%20C++\%20and\%20Systems\%20APIs\#Eigenvectors\%20reveal\%20the\%20directions\%20a\%20square\%20map\%20does\%20not\%20turn}{eigenanalysis}. - \hyperref[appendix-e8ab9836882e]{ODE initial-value problems and stiffness}

\subsection{Full teaching explanation}\label{full-teaching-explanation}

\begin{itemize}
\tightlist
\item
  \href{Continuous\%20Notes/Computing\%20Foundations\%20-\%20From\%20Bits\%20to\%20C++\%20and\%20Systems\%20APIs\#ODEs\%20and\%20numerical\%20simulation}{Phase 9 teaching section}
\item
  \href{}{Visual Recall/09 - Calculus and Optimization - Visual Recall.md (01-Foundations-I, appendix-adf48f5dc5c1)} · Phase 9 visual recall
\end{itemize}
